
\section{La réplication, pourquoi?}

\begin{frame}
\frametitle{Répliquer, à quoi ça sert\,?}

Outil indispensable, universel pour la robustesse des systèmes distribués.

\begin{redbullet}
  \item \alert{Tolérance aux pannes}
  
     Vous cherchez un document sur $S_1$, qui est en panne\,? On le trouvera sur $S_2$
     
 \item \alert{Distribution des lectures}
 
    Répartissons les lectures sur $S_1, S_2, \ldots, S_n$ pour satisfaire les millions de requêtes de nos clients.
    
 \item \alert{Distribution des écritures\,?}
 
    Oui, mais attention, il faut \alert{réconcilier} les données ensuite.
    
     \item et autres avantages, comme la construction d'un index sur un des serveurs sans affecter les autres, 
     
 \end{redbullet}

Le bon niveau de réplication? \alert{Trois copies pour une sécurité totale; deux au minimum}.

\end{frame}


\begin{frame}
\frametitle{Méthode générale de réplication}

Une \alert{application} (le client) demande au \alert{système} (le serveur) $S_1$ \alert{l'écriture}
d'un document (unité d'information).

\begin{redbullet}
  \item le serveur écrit le document sur le disque\,; 
  \item \alert{$S_1$ transmet la demande d'écriture à un ou plusieurs autres serveurs $S_2, \ldots, S_n$}, créant
     des \alert{replicas} ou \alert{copies}\,;
  \item $S_1$ ``rend la main'' au client.
\end{redbullet}

Essentiel, réplication synchrone/asynchrone\,?

\begin{redbullet}
  \item \alert{synchrone}: $S_1$ \alert{attend} confirmation de $S_2, \ldots, S_n$ avant de rendre la main\,; 
 \item \alert{asynchrone}: $S_1$ rend la main \alert{sans attendre}. 
 \end{redbullet}

\begin{remblock}{Conséquences}
Une réplication \alert{asynchrone} est beaucoup plus rapide, mais elle favorise les \alert{incohérences}, au moins
temporaires.
\end{remblock}


\end{frame}


\begin{frame}
\frametitle{Pour éviter la latence disque: le fichier journal}

Technique universellement utilisée en centralisé

\vfill


\figSlideWithSize{log}{11cm}

\vfill

Permet de se ramener à des entrées/sorties \alert{séquentielles}.

\end{frame}


\begin{frame}
\frametitle{Réplication avec écritures synchrones}

Le client est acquitté quand \alert{tous} les serveurs ont effectué l'écriture.

\vfill

\figSlideWithSize{replication-synchrone}{10cm}

\vfill

Sécurité totale; cohérence forte; \alert{très lent}.

\end{frame}

\begin{frame}
\frametitle{Réplication avec écritures asynchrones}

Le client est acquitté quand \alert{un} des serveurs a effectué l'écriture. Les autres
écritures se font indépendamment.

\vfill

\figSlideWithSize{replication-asynchrone}{10cm}

\vfill

Sécurité partielle; cohérence faible; \alert{Efficace}.

\end{frame}

\section{Cohérence des données}


\begin{frame}
\frametitle{Parlons cohérence}

La cohérence est la capacité d’un système de gestion de données à refléter fidèlement les opérations d’une application.
 
\begin{redbullet}
  \item toute opération (validée) est immédiatement visible et permanente. 
  \item si je fais une écriture de $d$ suivie d’une lecture, je dois constater les modifications effectuées; 
  \item si je refais une lecture un peu plus tard, ces modifications doivent toujours être présentes.
\end{redbullet}

\alert{Cohérence forte = } une des fonctionnalités marquantes des systèmes relationnels.

\begin{remblock}{Dans les systèmes NoSQL}
\alert{La cohérence forte est sacrifiée au profit des performances.} 
\end{remblock}
\end{frame}

\begin{frame}
\frametitle{Architectures avec réplication: topologie et (a)synchronicité}

\figSlideWithSize{coherence}{11cm}

\end{frame}


\begin{frame}
\frametitle{Soyons cohérents (ou pas...)}

Plusieurs niveaux de cohérence dans les systèmes distribués / NoSQL.

\vfill

\begin{redbullet}
  \item \alert{Cohérence forte} (ACID) --  nécessite une réplication synchrone,
     et potentiellement des verrouillages complexes (\textit{two phases commit}).
       
  \item \alert{Cohérence faible} -- on accepte le risque de lectures ne reflétant pas les mises à jour.
  
  \item \alert{Cohérence à terme} (\textit{Eventual consistency}) -- le système
     garantit que les incohérences ne sont que transitoires. 
\end{redbullet}

\begin{remblock}{Le NoSQL}
Les systèmes NoSQL dans l'ensemble abandonnent la cohérence forte pour
favoriser la réplication asynchrone, et donc le débit en écriture/lecture.
\end{remblock}

\end{frame}

\section{Reprise sur panne}

\begin{frame}
\frametitle{Réplication et reprise sur panne}

Principe général: tout le monde est interconnecté et échange des messages
(\textit{heartbeat}.

\figSlideWithSize{heartbeat}{10cm}

\vfill

  Si un \alert{esclave} tombe en panne, le système continue à fonctionner, \alert{tant qu'il est en contact
    avec une majorité d'esclaves}
 
 \vfill
 
 Si le \alert{maître} tombe en panne, les \alert{esclaves} doivent \alert{élire} 
 un nouveau maître.

\end{frame}

\begin{frame}{Cas du partitionnement réseau}

Un \alert{partitionnement} divise la grappe en deux groupes.

\figSlideWithSize{reprise}{10cm}

\begin{remblock}{Principe de majorité}
Seul le groupe représentant la \alert{majorité} des participants initiaux peut élire
un maître.
\end{remblock}

Election: algorithme de négociation, voir par exemple Paxos.

\end{frame}



\begin{frame}
\frametitle{Culture, culture\,: le théorème CAP}

Un ``théorème'' qui dit qu'un système distribué ne peut satisfaire à la fois

\begin{redbullet}
  \item La cohérence (\alert{C})
  
  \item La disponibilité (\alert{A}\textit{vailability})
   \item La tolérance au  \alert{P}artitionnement 
  
\end{redbullet}

Exprime une limitation intrinsèque aux systèmes distribués, applicable à tout système
NoSQL, et expliquant qu'il faut toujours faire un compromis.

\end{frame}


\end{document}


